\section{Exact certificates and limitations of local reductions}
\label{root:certificates}

The companion code uses integer arithmetic and the Python standard
library. Its certificates are finite statements with explicit
hypotheses. They do not convert an unsuccessful search into a
counterexample or a universal theorem.

\subsection{A unified sufficient recognizer}

The program \texttt{exact2n\_recognition.py} keeps the ambient
length \(M\) fixed. At a current set of \(m\) sources it recognizes
the full-interval base, Wang's external quadratic base, or
Theorem~\ref{root:hall-theorem}. Otherwise it may delete a source by
one of the following proved extension rules:
\[
 \Omega(a)\le2;
\]
\[
 pq\mid a,\quad pq\ge2m-1,\quad
 D_p(A)\le q\,v_p(a/q);
\]
\[
 a=p^e q,\quad p,q\text{ prime},\quad mp^e\le M.
\]
Each valid deletion is restored with at most two additional
positions, so a successful trace proves the exact \(2|A|\) bound
for every shift. The certificate verifier starts with the original
source list and recomputes demands, cofactors, numerical witnesses,
and all Hall rows. It does not rerun the recognizer's choice heuristic;
it does share the exact arithmetic helpers. If the rules stop, the
output is \texttt{UNKNOWN}. No cover is explicitly constructed by
this recognizer.

For a concrete combination example, let \(\mathcal R\) be all 35
cubic monomials, with repetitions, on
\(\{11,13,17,19,23\}\), and set
\[
 A=\mathcal R\cup\{4096,11875\},\qquad
 n=37,\quad M=12167.
\]
Every source has at least three prime factors, counting multiplicity.
The direct Hall certificate fails: at \(d=121\) its lower bound is
\(-7\), while the fibre has size five. Nevertheless,
\[
 11875=5^4\cdot19,\qquad
 5\cdot19=95\ge73,\qquad D_5(A)=4\le19\cdot4.
\]
The two-prime rule deletes this source. All Hall rows of the remaining
36 sources then pass; their minimum margin is zero. Hence every
length-12167 interval has a cover of size at most 74 for this set.
Wang's quadratic base is inapplicable both before and after deletion,
and the prime-power-bank rule alone cannot remove this source because
\(37\cdot625>12167\).

The saved examples also include deleting the maximum while retaining
the original ambient length, two successive deletion steps, direct
base cases, and an unresolved 64-source smooth input. Their validation
checks the combined proof mechanism rather than sampling interval
shifts.

\subsection{Necessary arithmetic cores}

The program \texttt{necessary\_core.py} applies necessary conditions
for a \emph{minimal-cardinality} counterexample. It computes exact
minimum forced-union cardinalities over a period and rigorous lower
bounds for their valuations. The latter are sums of individual level
minima; they are not asserted to equal the minimum total valuation at
one common shift. It then removes candidates whose required demand
exceeds the largest possible demand of a source set of the stated
cardinality.

For example, at
\[
 n=100,\qquad M=146969=47\cdot53\cdot59,
\]
the forced 47-valuation of the maximum source is at least 112, while
the largest possible 47-demand of 100 sources from the candidate pool,
including the maximum, is at most 90. The accompanying exact trace
therefore excludes these \emph{minimal-counterexample parameters}.
It does not establish \(g(100)\le200\).

There is also an exact source-only family showing the limits of
these filters. Its certificate has
\[
 n=16384,\qquad M=n^2=268435456,
\]
and uses all 30 primes at most 113. Its source and cofactor prime
supports coincide, its common gcd is one, and its prime-support
hypergraph is connected. Every source has at least three prime
occurrences, and every prime-deletion cofactor exceeds \(M/n\).
All prime-pair products are below \(2n-1\), so the two-prime demand
tests are vacuous. The forced-union cardinality bound is at most 676,
well below the available 32766 positions, while the cofactor
reciprocal mass exceeds one and the finite Hall certificate fails.

This is a witness to the insufficiency of the listed arithmetic
conditions. It is not a counterexample to the covering conjecture:
in particular at shift zero the original sources cover themselves.
No collection of deletion covers satisfying every universal
minimality requirement is asserted to exist.

\subsection{Why unrestricted cofactor matching cannot be the missing step}

Even for sources with at least three prime factors, the particular
assignment \(a\mapsto a/P^+(a)\) need not admit a matching. An exact
example retained from the preceding work is
\[
 M=440,\quad I=\{4401,\ldots,4840\},\quad
 A=\{2k:k\le220\text{ composite},\ k\ne169\}.
\]
There are 171 sources. The union of their cofactor neighborhoods has
only 170 positions: exactly the interval integers divisible by at
least one of \(4,6,10,14,22\). Thus Hall fails for this proposed
assignment, although the original product-cover problem allows
other constructions and a much larger budget. This obstruction
concerns the proposed matching, not the \(2n\) conjecture.

\section{The remaining mathematical target}
\label{root:remaining}

The exact and asymptotic goals must be kept distinct. The structural
core defined by source demands is not automatically the same set as
the large coordinates of an optimizing dual vector. Nor can separate
covers for two source parts be combined by adding their valuation
contributions: an overlapping target is only available once.

\subsection{An exact criterion for asymptotic closure}

The raw boxed dual in Lemma~\ref{an:lp-bridge}, together with Wang's
uniform raw fractional transfer theorem \cite{Wang}, gives the
following useful precise target, also recorded in \cite{Earlier}.

\begin{proposition}[Uniform additive-defect criterion]
\label{root:closure-criterion}
The unrestricted assertion \(g(n)\le(2+o(1))n\) is equivalent to
\begin{equation}
 \forall\varepsilon>0\ \exists C_\varepsilon<\infty\
 \forall M,I,w:\qquad
 \Delta_{2+\varepsilon}(M,I,w)\le C_\varepsilon,
 \label{root:closure-target}
\end{equation}
where \(w\) has finite support and \(0\le w_p\le1\), and
\(C_\varepsilon\) depends only on \(\varepsilon\).
\end{proposition}

\begin{proof}
The asymptotic assertion implies, for each \(\varepsilon>0\), a
bound \(g(n)\le(2+\varepsilon)n+C_\varepsilon\) for all \(n\),
absorbing finitely many small cardinalities into the constant. It
bounds every boxed dual objective for every
\(A\subseteq[2,M]\). If \(\max A<M\), use a contained interval
of length \(\max A\), whose available target hinge sum is at most
that on \(I\). Maximize over all such subsets and use
\[
 \max_A\left\{\sum_{a\in A}X_w(a)-(2+\varepsilon)|A|\right\}
 =\sum_{a\le M}(X_w(a)-2-\varepsilon)_+.
\]
This proves \eqref{root:closure-target}.

Conversely, \eqref{root:closure-target} and the boxed dual give
\(L(A,I)\le(2+\varepsilon)n+C_\varepsilon\) uniformly. For
completeness, the specific external transfer inequality being used
is as follows. Put \(\ell(0)=0\), and let \(\ell(m)\) be the
supremum of the raw boxed values on original \(m\)-source
instances. Then
\[
 K_n=\left\lceil\max_{0\le m\le n}
        \{\ell(m)+2(n-m)\}\right\rceil,
 \qquad g(n)\le K_n+\pi(2K_n)-1.
\]
Since \(\ell(m)\le(2+\varepsilon)m+C_\varepsilon\), this yields
\(g(n)\le(2+\varepsilon)n+C_\varepsilon+o(n)\).
Let \(\varepsilon\) tend to zero.
\end{proof}

\subsection{Rounding is already small on the reduced obstruction class}

For the minimal asymptotic-obstruction sequences constructed above,
all source primes are below \(K<10n\). Hence
\[
 |P_A|\le\pi(10n)=O(n/\log n),\qquad
 \tau(A,I)\le L(A,I)+O(n/\log n),
\]
by Lemma~\ref{an:lp-bridge}. This conclusion concerns these reduced
instances; it does not assert small rounding loss for an arbitrary
original instance.

Consequently a more focused sufficient next theorem is a uniform
fractional bound
\[
 L(A,I)\le(2+\delta)n+o(n)
\]
on this explicit structural class, for every fixed \(\delta>0\).
The error must be uniform in the source set, shift, and additive
minimality parameter. Indeed, use the minimality budget
\((2+2\delta)n+C\), with \(2+2\delta<10\). As \(C\to\infty\),
the source cardinality tends to infinity, and the displayed estimates
would give \(\tau(A,I)<(2+2\delta)n\), a contradiction. Thus the
outstanding issue on this class is its joint fractional valuation
bound. That bound is not proved here.

\subsection{A concrete order of attack}

The following are mathematical milestones, rather than a claim
that the remaining work is routine.

\begin{enumerate}
\item \textbf{Range and proof audit: completed.} The submitted
\(n^{5/2}\) proof is reconstructed, its essential restriction is
identified, and the independently checked quadratic input sets the
remaining range to \eqref{root:range}.
\item \textbf{Arbitrarily many small weights: completed.} The
diffuse theorem controls this class at every scale, including one
arbitrary exceptional coordinate. No bound on the number of prime
factors is assumed.
\item \textbf{A controlled large-weight block: completed within
its stated range.} The explicit defect allows a fixed or sufficiently
slowly growing number of exceptions. It does not handle a block as
large as the source-demand core in the combinatorial reduction.
\item \textbf{Dense small-prime sources and separated bands:
completed within their stated classes.} The smooth-source theorem
and the exact separated-band theorem are actual covering results.
Neither covers arbitrary combinations of adjacent prime scales.
\item \textbf{Growing mixed cores: open.} The next substantive
target is a uniform certificate handling substantial weights on
intermediate prime scales while retaining a growing dense small-prime
part. A uniform bound \eqref{root:closure-target} would close the
asymptotic question. An exact \(2n\) construction would additionally
have to control the integer position budget without an asymptotic
rounding allowance.
\end{enumerate}

One quantitative boundary of the present proof is explicit. In the
smooth-support argument, the sieve and tail errors behave like
\(\exp(-c\log M/\log Y)\), whereas the guaranteed model gap is
\(Y^{-J_\varepsilon}\). Their comparison works when
\((\log Y)^2\ll\log M\). It does not work at a general polynomial
cutoff \(Y=M^\alpha\). The separated-band argument handles one
part of the interaction by producing disjoint exact-part bins, but
it requires a large gap between the two prime ranges and one large
occurrence per source.

The package therefore contains verified partial theorems and a
precise remaining target, not a proof that problem 708 is solved or
that a numerical fraction of it has been completed.
